\chapter*{Anexos} \label{anex}
\addcontentsline{toc}{chapter}{Anexos}  % Añadir los anexos al índice sin numerarlos

Cambios de la gramática de C\#.

\section*{C\# 7.0}

\subsection*{Variables \texttt{out}}
\textbf{Problema:} El uso de variables \texttt{out} requería una declaración previa. Desde C\# 7.0, se permite declarar variables directamente en la declaración del método.

\textbf{Solución:} Se modificó la regla \texttt{argument} para incluir una declaración de variable opcional antes del modificador \texttt{out}. \\
 Además, se ajustó la regla \texttt{local\_variable\_declaration} para permitir el uso de \texttt{out} como modificador de variable.

\subsection*{Tuplas y Deconstrucción}
\textbf{Problema:} Las tuplas no eran un tipo nativo y requerían una biblioteca externa. Desde C\# 7.0, se introdujeron tuplas como un tipo nativo, permitiendo la deconstrucción de sus elementos.

\textbf{Solución:} Se agregó una nueva regla \texttt{tuple\_type} para definir tipos de tuplas. Se modificó la regla \texttt{type} para incluir \texttt{tuple\_type}. Además, se añadió una regla \texttt{tuple\_literal} para la creación de tuplas y se actualizó la regla \texttt{assignment} para soportar la deconstrucción de tuplas.

\subsection*{Coincidencia de Patrones}
\textbf{Problema:} No había soporte para coincidencia de patrones en C\#. Desde C\# 7.0, se introdujo la coincidencia de patrones que permite simplificar el código al verificar tipos y valores.

\textbf{Solución:} Se introdujo una nueva regla \texttt{pattern\_matching\_expression} que incluye subreglas para diferentes tipos de patrones (constante, tipo, var). Se actualizaron las reglas \texttt{is\_expression} y \texttt{switch\_section} para incorporar la nueva sintaxis de coincidencia de patrones.

\subsection*{Funciones Locales}
\textbf{Problema:} Las funciones debían estar definidas a nivel de clase. Desde C\# 7.0, se introdujeron funciones locales que permiten definir funciones dentro de métodos.

\textbf{Solución:} Se creó una nueva regla \texttt{local\_function\_declaration} similar a \texttt{method\_declaration}, pero permitida dentro de bloques de código. Se actualizó la regla \texttt{statement} para incluir \texttt{local\_function\_declaration} como una opción válida.

\subsection*{Miembros Expresados con Cuerpo de Expresión Expandido}
\textbf{Problema:} Los miembros debían tener un cuerpo completo. Desde C\# 7.0, se introdujo la posibilidad de usar cuerpos de expresión para simplificar la definición de miembros.

\textbf{Solución:} Se modificaron las reglas \texttt{method\_declaration}, \texttt{property\_declaration}, \texttt{operator\_declaration}, \texttt{constructor\_declaration}, \texttt{destructor\_declaration}, y \texttt{indexer\_declaration} para permitir el uso de \texttt{=>} seguido de una expresión como alternativa al cuerpo del bloque tradicional.

\subsection*{Locales \texttt{ref}}
\textbf{Problema:} No era posible utilizar locales \texttt{ref} que permiten referenciar variables directamente en lugar de copiar su valor. Desde C\# 7.0, se introdujeron los locales \texttt{ref}.

\textbf{Solución:} Se actualizó la regla \texttt{local\_variable\_declaration} para incluir el modificador \texttt{ref} opcional. Se agregó una nueva regla \texttt{ref\_expression} para manejar expresiones que devuelven referencias.

\subsection*{Retornos \texttt{ref}}
\textbf{Problema:} Los métodos no podían devolver referencias directamente. Desde C\# 7.0, se introdujo el soporte para retornos \texttt{ref}, permitiendo que un método devuelva una referencia a un valor existente.

\textbf{Solución:} Se modificó la regla \texttt{method\_declaration} para permitir el modificador \texttt{ref} antes del tipo de retorno. Se actualizó la regla \texttt{return\_statement} para incluir la opción de devolver una referencia usando \texttt{ref}.

\subsection*{Discards}
\textbf{Problema:} No había una forma explícita de ignorar valores en asignaciones o expresiones. Desde C\# 7.0, se introdujo el uso de discards (\texttt{\_}) que permite omitir valores que no son necesarios.

\textbf{Solución:} Se agregó una nueva regla \texttt{discard} que reconoce el símbolo \texttt{\_}. Se actualizaron las reglas relevantes como \texttt{variable\_declaration}, \texttt{foreach\_statement}, y \texttt{out\_variable\_declaration} para permitir el uso de \texttt{discard} en lugar de identificadores.

\subsection*{Literales Binarios y Separadores de Dígitos}
\textbf{Problema:} No había soporte para literales binarios ni separadores de dígitos en números. Desde C\# 7.0, se introdujeron literales binarios y la capacidad de usar guiones bajos como separadores en literales numéricos.

\textbf{Solución:} Se modificó la regla \texttt{integer\_literal} para incluir el prefijo \texttt{0b} o \texttt{0B} para literales binarios. Se actualizaron las reglas de literales numéricos para permitir guiones bajos (\texttt{\_}) entre dígitos en cualquier posición excepto al principio o final del literal.

\subsection*{Expresiones \texttt{throw}}
\textbf{Problema:} Las expresiones \texttt{throw} solo podían usarse en instrucciones completas. Desde C\# 7.0, se introdujeron expresiones \texttt{throw}, permitiendo lanzar excepciones como parte de expresiones más complejas.

\textbf{Solución:} Se creó una nueva regla \texttt{throw\_expression} que permite el uso de \texttt{throw} seguido de una expresión. Se actualizaron las reglas relevantes como \texttt{conditional\_expression}  y  \texttt{null\_coalescing\_expression}  para  incluir \\
\texttt{throw\_expression} como una opción válida.


\subsection*{C\# 7.1}

\subsection*{Método \texttt{Main} Asíncrono}
\textbf{Problema:} El método de entrada \texttt{Main} no podía ser asíncrono. Esto limitaba el uso de operaciones asincrónicas en la inicialización de aplicaciones. Desde C\# 7.1, se permite que el método \texttt{Main} tenga el modificador \texttt{async}.

\textbf{Solución:} Se ajustaron las reglas gramaticales para permitir el modificador \texttt{async} en el método de entrada \texttt{Main}.

\subsection*{Expresiones Literales por Defecto (\texttt{default})}
\textbf{Problema:} En versiones anteriores, era necesario especificar explícitamente el tipo al usar valores predeterminados. Desde C\# 7.1, se introdujeron las expresiones literales \texttt{default}, que infieren automáticamente el tipo cuando es posible.

\textbf{Solución:} Se modificó la gramática para permitir expresiones literales \texttt{default} en contextos donde el tipo puede inferirse.

\subsection*{Inferencia de Nombres de Elementos en Tuplas}
\textbf{Problema:} En versiones anteriores, era necesario nombrar explícitamente los elementos de una tupla al inicializarla. Desde C\# 7.1, se permite inferir los nombres de los elementos a partir de las variables utilizadas en la inicialización.

\textbf{Solución:} Se ajustaron las reglas gramaticales para inferir automáticamente los nombres de los elementos en inicializaciones de tuplas.

\subsection*{Coincidencia de Patrones en Parámetros Genéricos}
\textbf{Problema:} No era posible realizar coincidencia de patrones en variables cuyo tipo era un parámetro genérico. Desde C\# 7.1, se permite realizar coincidencias de patrones en este caso.

\textbf{Solución:} Se expandieron las reglas para permitir coincidencia de patrones en variables con tipos genéricos.


\subsection*{C\# 7.2}

\subsection*{Inicializadores en Arreglos Stackalloc}
\textbf{Problema:} No era posible inicializar directamente los arreglos creados con \texttt{stackalloc}. Desde C\# 7.2, se permite usar inicializadores en arreglos \texttt{stackalloc}.

\textbf{Solución:} Se ajustó la gramática para permitir inicializadores en arreglos creados con \texttt{stackalloc}.

\subsection*{Uso de Sentencias Fijas con Tipos que Soportan Patrones}
\textbf{Problema:} Las sentencias \texttt{fixed} solo podían usarse con tipos específicos como arreglos o cadenas. Desde C\# 7.2, se permite usar sentencias \texttt{fixed} con cualquier tipo que soporte un patrón.

\textbf{Solución:} Se modificó la regla \texttt{fixed\_statement} para permitir el uso de cualquier tipo que soporte patrones.

\subsection*{Acceso a Campos Fijos Sin Fijación}
\textbf{Problema:} Para acceder a campos fijos era necesario fijar la dirección de memoria explícitamente. Desde C\# 7.2, se permite acceder a campos fijos sin necesidad de fijarlos.

\textbf{Solución:} Se ajustaron las reglas gramaticales para permitir el acceso a campos fijos sin fijación explícita.

\subsection*{Reasignación de Variables Locales Referenciadas}
\textbf{Problema:} No era posible reasignar variables locales referenciadas (\texttt{ref}). Desde C\# 7.2, se permite reasignar estas variables.

\textbf{Solución:} Se modificaron las reglas para permitir la reasignación de variables locales referenciadas.

\subsection*{Declaración de Tipos Struct Solo de Lectura (\texttt{readonly struct})}
\textbf{Problema:} No había una forma explícita de declarar estructuras inmutables. Desde C\# 7.2, se introdujo el modificador \texttt{readonly struct} para indicar que una estructura es inmutable y debe pasarse como un parámetro \texttt{in}.

\textbf{Solución:} Se agregó soporte para el modificador \texttt{readonly} en la declaración de estructuras.

\subsection*{Modificador \texttt{in} en Parámetros}
\textbf{Problema:} No era posible pasar argumentos por referencia sin permitir su modificación. Desde C\# 7.2, se introdujo el modificador \texttt{in}, que permite pasar argumentos por referencia sin permitir modificaciones.

\textbf{Solución:} Se ajustó la regla \texttt{parameter\_modifier} para incluir el modificador \texttt{in}.

\subsection*{Modificador \texttt{ref readonly} en Retornos de Métodos}
\textbf{Problema:} No había una forma explícita de devolver referencias inmutables desde métodos. Desde C\# 7.2, se introdujo el modificador \texttt{ref readonly}, que permite devolver referencias inmutables.

\textbf{Solución:} Se agregó soporte para el modificador \texttt{ref readonly} en las reglas de retorno de métodos.

\subsection*{Declaración de Tipos Struct Referenciados (\texttt{ref struct})}
\textbf{Problema:} No había una forma explícita de declarar estructuras que accedan directamente a memoria administrada y deban ser asignadas en el stack. Desde C\# 7.2, se introdujo el modificador \texttt{ref struct}.

\textbf{Solución:} Se agregó soporte para el modificador \texttt{ref struct} en la declaración de estructuras.

\subsection*{Argumentos Nombrados No Finales}
\textbf{Problema:} En versiones anteriores, los argumentos nombrados debían colocarse al final de la lista de argumentos. Desde C\# 7.2, se permite mezclar argumentos posicionales y nombrados siempre que los posicionales estén primero.

\textbf{Solución:} Se ajustaron las reglas gramaticales para permitir argumentos nombrados no finales.

\subsection*{Guiones Bajos Iniciales en Literales Numéricos}
\textbf{Problema:} Los literales numéricos no podían comenzar con guiones bajos (\texttt{\_}). Desde C\# 7.2, se permite usar guiones bajos iniciales antes de los dígitos impresos.

\textbf{Solución:} Se ajustó la gramática para permitir guiones bajos iniciales en literales numéricos.

\subsection*{Modificador \texttt{private protected}}
\textbf{Problema:} No había un modificador que combinara las restricciones de acceso \texttt{private} y \texttt{protected}. Desde C\# 7.2, se introdujo el modificador \texttt{private protected}, que permite acceso desde clases derivadas dentro del mismo ensamblado.

\textbf{Solución:} Se agregó soporte para el modificador \texttt{private protected} en las reglas de acceso.

\subsection*{Expresiones Condicionales por Referencia}
\textbf{Problema:} En versiones anteriores, las expresiones condicionales (\texttt{?:}) no podían devolver referencias. Desde C\# 7.2, se permite que el resultado sea una referencia.

\textbf{Solución:} Se ajustó la gramática para permitir expresiones condicionales por referencia.


\subsection*{C\# 7.3}

\subsection*{Acceso a Campos Fijos Sin Fijación}
\textbf{Problema:} En versiones anteriores de C\#, para acceder a campos fijos era necesario fijar la dirección de memoria. Desde C\# 7.3, se permite acceder a campos fijos sin necesidad de fijarlos explícitamente.

\textbf{Solución:} Se ajustaron las reglas gramaticales para permitir el acceso a campos fijos sin fijación.

\subsection*{Reasignación de Variables Locales Referenciadas}
\textbf{Problema:} No se podía reasignar variables locales referenciadas. Desde C\# 7.3, se permite la reasignación de variables locales que están marcadas como \texttt{ref}.

\textbf{Solución:} Se modificaron las reglas para permitir la reasignación de variables locales referenciadas.

\subsection*{Inicializadores en Arreglos Stackalloc}
\textbf{Problema:} No era posible utilizar inicializadores con arreglos creados con \texttt{stackalloc}. Desde C\# 7.3, se permite usar inicializadores en arreglos \texttt{stackalloc}.

\textbf{Solución:} Se ajustaron las reglas gramaticales para permitir inicializadores en arreglos \texttt{stackalloc}.

\subsection*{Sentencias Fijas con Tipos que Soportan Patrones}
\textbf{Problema:} Las sentencias fijas solo podían usarse con tipos específicos. Desde C\# 7.3, se permiten sentencias fijas con cualquier tipo que soporte patrones.

\textbf{Solución:} Se modificó la regla \texttt{fixed\_statement} para permitir el uso con cualquier tipo que soporte patrones.

\subsection*{Pruebas con Tipos Tupla}
\textbf{Problema:} No se podía usar \texttt{==} y \texttt{!=} directamente con tipos tupla. Desde C\# 7.3, se permite comparar tuplas usando estos operadores.

\textbf{Solución:} Se modificaron las reglas para permitir comparaciones directas entre tipos tupla.

\subsection*{Uso de Variables de Expresión en Más Ubicaciones}
\textbf{Problema:} Las variables de expresión tenían un alcance limitado. Desde C\# 7.3, se permite usar variables de expresión en más ubicaciones dentro del código.

\textbf{Solución:} Se ajustaron las reglas gramaticales para ampliar el uso de variables de expresión en diferentes contextos.

\subsection*{Resolución de Métodos Mejorada al Usar Argumentos Diferentes por \texttt{in}}
\textbf{Problema:} La resolución de métodos al usar argumentos diferentes por \texttt{in} era confusa y poco clara. Desde C\# 7.3, se mejoró este aspecto.

\textbf{Solución:} Se ajustaron las reglas gramaticales para mejorar la resolución de métodos al usar argumentos por \texttt{in}.

\section*{C\# 8}

\subsection*{Miembros Solo de Lectura \texttt{readonly}}
\textbf{Problema:} No se podía definir miembros de solo lectura en las clases. Desde C\# 8, se introdujeron los miembros solo de lectura, que permiten que las propiedades sean inmutables después de su inicialización.

\textbf{Solución:} Se modificó la regla \texttt{property\_declaration} para permitir miembros solo de lectura \texttt{readonly}.

\subsection*{Métodos de Interfaz por Defecto}
\textbf{Problema:} No era posible definir métodos en interfaces con una implementación por defecto. Desde C\# 8, se introdujeron métodos de interfaz por defecto, permitiendo que las interfaces contengan métodos implementados.

\textbf{Solución:} Se ajustó la regla \texttt{interface\_declaration} para permitir métodos de interfaz con implementación por defecto.

\subsection*{Mejoras en Coincidencia de Patrones}
\textbf{Problema:} Las versiones anteriores de C\# tenían limitaciones en los patrones disponibles. Desde C\# 8, se introdujeron mejoras en la coincidencia de patrones, incluyendo expresiones \texttt{switch} y patrones de propiedades.

\textbf{Solución:} Se expandieron las reglas \texttt{pattern} para incluir expresiones \texttt{switch} y patrones mejorados.

\subsection*{Patrones de Propiedades}
\textbf{Problema:} No había soporte para patrones que coincidieran con propiedades específicas. Desde C\# 8, se introdujeron patrones de propiedades para facilitar la coincidencia con las propiedades de un objeto.

\textbf{Solución:} Se modificó la regla \texttt{pattern} para incluir patrones que coincidan con propiedades.

\subsection*{Patrones Tupla}
\textbf{Problema:} No había soporte para coincidencias basadas en tuplas. Desde C\# 8, se introdujeron patrones tupla para facilitar la coincidencia con estructuras de tuplas.

\textbf{Solución:} Se agregó soporte para patrones tupla en la regla \texttt{pattern}.

\subsection*{Patrones Posicionales}
\textbf{Problema:} No había soporte para patrones que coincidieran con la posición de los elementos. Desde C\# 8, se introdujeron patrones posicionales que permiten coincidir elementos en una posición específica.

\textbf{Solución:} Se agregó soporte para patrones posicionales en la regla \texttt{pattern}.

\subsection*{Declaraciones Using}
\textbf{Problema:} Las declaraciones \texttt{using} debían estar al principio del archivo o dentro del bloque. Desde C\# 8, se permite utilizar declaraciones \texttt{using} dentro del bloque de código.

\textbf{Solución:} Se modificó la regla \texttt{using\_declaration} para permitir declaraciones \texttt{using} más flexibles.

\subsection*{Funciones Locales Estáticas}
\textbf{Problema:} Las funciones locales no podían ser estáticas antes. Desde C\# 8, se permite definir funciones locales como estáticas.

\textbf{Solución:} Se agregó soporte para funciones locales estáticas en la gramática.

\subsection*{Estructuras Ref Descartables}
\textbf{Problema:} No había un manejo específico para estructuras ref descartables. Desde C\# 8, se introdujo el soporte para estructuras ref descartables que permiten un manejo más eficiente de la memoria.

\textbf{Solución:} Se modificó la regla \texttt{struct\_declaration} para permitir estructuras ref descartables.

\subsection*{Tipos Nullable Referencia}
\textbf{Problema:} Todas las referencias eran consideradas no nulas. Desde C\# 8, se introdujo el concepto de tipos nullable referencia que permite indicar explícitamente si una referencia puede ser nula.

\textbf{Solución:} Se ajustaron las reglas gramaticales para incluir el manejo de tipos nullable referencia.

\subsection*{Flujos Asincrónicos}
\textbf{Problema:} No había un manejo eficiente para flujos asincrónicos. Desde C\# 8, se introdujo el soporte para flujos asincrónicos que permite trabajar con secuencias de datos asincrónicas.

\textbf{Solución:} Se modificaron las reglas gramaticales para permitir flujos asincrónicos y su uso en bucles.

\subsection*{Índices y Rangos}
\textbf{Problema:} No había una forma sencilla de trabajar con índices y rangos en colecciones. Desde C\# 8, se introdujo el soporte para índices y rangos que facilita el acceso a elementos en colecciones.

\textbf{Solución:} Se agregó soporte para índices y rangos en la gramática.

\subsection*{Asignación Null-Coalescing}
\textbf{Problema:} No había forma concisa de asignar valores cuando una variable es nula. Desde C\# 8, se introdujo el operador null-coalescing assignment (\texttt{??=}).

\textbf{Solución:} Se ajustó la gramática para incluir el operador null-coalescing assignment.

\subsection*{Tipos Construidos No Administrados}
\textbf{Problema:} No había un manejo específico para tipos construidos no administrados. Desde C\# 8, se introdujo el soporte para tipos construidos no administrados que permiten trabajar con tipos sin administración del recolector de basura.

\textbf{Solución:} Se agregó soporte para tipos construidos no administrados en la gramática.

\subsection*{Stackalloc en Expresiones Anidadas}
\textbf{Problema:} No era posible usar \texttt{stackalloc} dentro de expresiones anidadas. Desde C\# 8, se permite usar \texttt{stackalloc} dentro de expresiones anidadas.

\textbf{Solución:} Se ajustó la gramática para permitir el uso de \texttt{stackalloc} en expresiones anidadas.

\subsection*{Mejora en las Cadenas Interpoladas Verbatim}
\textbf{Problema:} Las cadenas interpoladas verbatim tenían limitaciones en su formato. Desde C\# 8, se mejoraron las cadenas interpoladas verbatim para permitir más flexibilidad al incluir expresiones complejas.

\textbf{Solución:} Se modificó la regla \texttt{interpolated\_string\_expression} para mejorar el manejo de cadenas interpoladas verbatim.



\section*{C\# 9}

\subsection*{Registros \texttt{record}}

\textbf{Problema:} No existía un tipo de datos específico para registros. Sin embargo, desde C\# 9, se introdujeron los registros, que son tipos de referencia inmutables que facilitan la creación de objetos con propiedades.

\textbf{Solución:} Se modificó la regla \texttt{class\_declaration} para permitir la definición de \texttt{record}.

\subsection*{Setters Solo de Inicialización \texttt{init}}

\textbf{Problema:} Todas las propiedades podían ser modificadas después de su inicialización. Desde C\# 9, se introdujo la posibilidad de definir propiedades con setters solo de inicialización, permitiendo que las propiedades sean asignadas solo durante la creación del objeto.

\textbf{Solución:} Se ajustó la regla \texttt{property\_declaration} para incluir el modificador \texttt{init}.

\subsection*{Declaraciones en Nivel Superior}

\textbf{Problema:} El código debía estar contenido dentro de una clase o espacio de nombres. Desde C\# 9, se permite escribir declaraciones en nivel superior, lo que simplifica el código.

\textbf{Solución:} Se modificó la gramática para permitir declaraciones en nivel superior sin necesidad de definir una clase.

\subsection*{Mejoras en Coincidencia de Patrones}

\textbf{Problema:} Las versiones anteriores de C\# tenían limitaciones en los patrones disponibles. Desde C\# 9, se introdujeron mejoras en la coincidencia de patrones, incluyendo patrones relacionales y lógicos.

\textbf{Solución:} Se expandieron las reglas \texttt{pattern} para incluir patrones relacionales (>, <, $\leq$, $\geq$) y lógicos(or, and, not).

\subsection*{Enteros de Tamaño Nativo}

\textbf{Problema:} No había un tipo específico para enteros que se adaptara al tamaño nativo del sistema. Desde C\# 9, se introdujeron enteros de tamaño nativo (\texttt{nint} y \texttt{nuint}).

\textbf{Solución:} Se agregó soporte para los tipos \texttt{nint} y \texttt{nuint} en la regla \texttt{type\_}.

\subsection*{Punteros a Funciones}

\textbf{Problema:} No era posible trabajar con punteros a funciones directamente. Desde C\# 9, se introdujo el soporte para punteros a funciones.

\textbf{Solución:} Se agregó soporte para punteros a funciones en las reglas correspondientes.

\subsection*{Inicializadores de Módulo}

\textbf{Problema:} No había una forma específica de inicializar módulos. Desde C\# 9, se introdujeron inicializadores de módulo.

\textbf{Solución:} Se agregó soporte para inicializadores de módulo \texttt{[ModuleInitializer]} en la gramática.

\subsection*{Nuevas Funciones para Métodos Parciales}

\textbf{Problema:} Las funciones parciales tenían limitaciones en su uso. Desde C\# 9, se introdujeron nuevas características para métodos parciales.

\textbf{Solución:} Se ajustaron las reglas relacionadas con métodos parciales para incluir nuevas características.

\subsection*{Expresiones Nuevas con Tipos Objetivo}

\textbf{Problema:} No había una forma concisa de crear instancias utilizando tipos objetivo. Desde esta versión, se permite crear expresiones tipadas directamente.

\textbf{Solución:} Se ajustó la gramática para permitir expresiones tipadas al usar \texttt{new}.

\subsection*{Funciones Anónimas Estáticas}

\textbf{Problema:} Las funciones anónimas no podían ser estáticas anteriormente. Sin embargo, desde C\# 9, se permite definir funciones anónimas como estáticas.

\textbf{Solución:} Se agregó soporte para funciones anónimas estáticas en la gramática.

\subsection*{Expresiones Condicionales Tipadas por Objetivo}

\textbf{Problema:} Las expresiones condicionales no podían ser tipadas por objetivo antes. Desde C\# 9, se permite que las expresiones condicionales tengan un tipo objetivo inferido.

\textbf{Solución:} Se modificó la gramática para permitir expresiones condicionales tipadas por objetivo.

\subsection*{Tipos Retornados Covariantes}

\textbf{Problema:} Los tipos retornados covariantes \( c \, ? \, e_1 \, : \, e_2 \)
 no eran compatibles. Sin embargo, desde C\# 9, se introdujo el soporte para tipos retornados covariantes en métodos virtuales o sobrescritos.

\textbf{Solución:} Se ajustó la regla \texttt{method\_declaration} para permitir tipos retornados covariantes.

\subsection*{Soporte para Extensiones GetEnumerator en Bucles \texttt{foreach}}

\textbf{Problema:} El bucle \texttt{foreach} requería que las colecciones implementaran explícitamente \texttt{IEnumerable}. Con esta versión, se permite un soporte más flexible mediante extensiones \texttt{GetEnumerator}.

\textbf{Solución:} Se modificó la regla correspondiente al bucle \texttt{foreach} para incluir soporte para extensiones \texttt{GetEnumerator}.

\subsection*{Parámetros Descartados en Lambdas}

\textbf{Problema:} Los parámetros descartados no eran compatibles en expresiones lambda. Sin embargo, desde C\# 9, se permite usar parámetros descartados (\texttt{\_}) en lambdas.

\textbf{Solución:} Se ajustó la regla \texttt{lambda\_expression} para permitir parámetros descartados.

\subsection*{Atributos en Funciones Locales}

\textbf{Problema:} Anteriormente no era posible aplicar atributos a funciones locales. Sin embargo, desde C\# 9, se permite aplicar atributos a funciones locales definidas dentro de métodos.

\textbf{Solución:} Se modificó la regla correspondiente a los atributos para permitir su uso en funciones locales.

\subsection*{Uso del operador \texttt{with}}
\textbf{Problema:} Antes de C\# 9.0, no existía una forma concisa y directa de crear copias de objetos inmutables con algunas propiedades modificadas. Esto dificultaba la manipulación de objetos inmutables, ya que requería la creación manual de nuevos objetos con las propiedades deseadas.

\textbf{Solución:} Se modificó la regla \texttt{object\_initializer} para permitir el uso del operador \texttt{with}, que permite a los desarrolladores crear un nuevo objeto basado en uno existente y modificar solo las propiedades especificadas. Esto se implementó mediante la adición de una nueva regla \texttt{with\_expression} en la gramática, que permite la sintaxis \texttt{objeto with \{ propiedad = valor \}}.

\section*{C\# 10}

\subsection*{Estructuras de Registro \texttt{record struct}}
\textbf{Problema:} Las estructuras de registro no estaban disponibles. Sin embargo, desde C\# 10, se introdujeron las estructuras de registro, que permiten definir estructuras con características de registro.

\textbf{Solución:} Se modificó la regla \texttt{struct\_declaration} para permitir estructuras de registro.

\subsection*{Manejadores de Cadenas Interpoladas}
\textbf{Problema:} No había una forma explícita de manejar las cadenas interpoladas. Sin embargo, desde C\# 10, se introdujeron los manejadores de cadenas interpoladas para mejorar su manejo.

\textbf{Solución:} Se ajustó la regla \texttt{interpolated\_string\_expression} para reconocer correctamente las expresiones dentro de cadenas interpoladas.

\subsection*{Directivas \texttt{using} Globales}
\textbf{Problema:} Las directivas \texttt{using} debían estar dentro de un espacio de nombres. Sin embargo, desde C\# 10, se permiten directivas \texttt{using} globales que pueden aplicarse a todo el archivo.

\textbf{Solución:} Se modificó la regla \texttt{using\_directive} para permitir directivas \texttt{using} globales.

\subsection*{Declaración de Espacios de Nombres \texttt{namespace} a Nivel de Archivo}
\textbf{Problema:} Los espacios de nombres debían declararse utilizando llaves (\texttt{\{\}}). Sin embargo, desde C\# 10, también es posible declararlos con punto y coma (\texttt{;}), permitiendo definiciones más concisas.

\textbf{Solución:} Se amplió la regla \texttt{namespace\_declaration} para soportar ambas sintaxis.

\subsection*{Patrones de Propiedades Extendidos}
\textbf{Problema:} Los patrones de propiedades no eran tan flexibles. Sin embargo, desde C\# 10, se extendieron los patrones de propiedades para permitir más flexibilidad en la coincidencia de patrones.

\textbf{Solución:} Se modificó la regla \texttt{pattern} para soportar patrones de propiedades extendidos \texttt{property\_pattern}.

\subsection*{Expresiones Lambda con Tipo Natural}
\textbf{Problema:} Las expresiones lambda debían tener un tipo explícito. Sin embargo, desde C\# 10, se permite que las expresiones lambda tengan un tipo natural, donde el compilador puede inferir el tipo delegado.

\textbf{Solución:} Se modificó la regla \texttt{lambda\_expression} para permitir tipos naturales en expresiones lambda.

\subsection*{Expresiones Lambda con Tipo de Retorno Explícito}
\textbf{Problema:} Las expresiones lambda no podían declarar un tipo de retorno explícito. Sin embargo, desde C\# 10, se permite que las expresiones lambda declaren un tipo de retorno cuando el compilador no puede inferirlo.

\textbf{Solución:} Se modificó la regla \texttt{lambda\_expression} para permitir tipos de retorno explícitos.

\subsection*{Atributos en Expresiones Lambda}
\textbf{Problema:} Los atributos no podían aplicarse a expresiones lambda. Sin embargo, desde C\# 10, se permite aplicar atributos a expresiones lambda.

\textbf{Solución:} Se modificó la regla \texttt{lambda\_expression} para permitir atributos.

\subsection*{Inicialización de Cadenas Constantes con Interpolación}
\textbf{Problema:} Las cadenas constantes no podían inicializarse con interpolación. Sin embargo, desde C\# 10, se permite inicializar cadenas constantes utilizando interpolación si todos los marcadores son cadenas constantes.

\textbf{Solución:} Se ajustó la regla \texttt{constant\_expression} para permitir interpolación en cadenas constantes.

\subsection*{Modificador \texttt{sealed} en \texttt{ToString} de Registros}
\textbf{Problema:} No había una forma explícita de sellar el método \texttt{ToString} en tipos de registro. Sin embargo, desde C\# 10, se permite agregar el modificador \texttt{sealed} cuando se sobrescribe \texttt{ToString} en un tipo de registro.

\textbf{Solución:} Se modificó la regla \texttt{method\_declaration} para permitir el modificador \texttt{sealed} en \texttt{ToString}.

\subsection*{Advertencias para Asignación Definida y Análisis de Estado Nulo}
\textbf{Problema:} Las advertencias para asignación definida y análisis de estado nulo no eran tan precisas. Sin embargo, desde C\# 10, se mejoraron estas advertencias para ser más precisas.

\textbf{Solución:} Se ajustó la regla \texttt{expression} para mejorar las advertencias relacionadas con la asignación definida y el análisis de estado nulo.

\subsection*{Declaración y Asignación en Deconstrucción}
\textbf{Problema:} La deconstrucción solo permitía declarar variables. Sin embargo, desde C\# 10, se permite tanto declarar como asignar valores en la misma deconstrucción.

\textbf{Solución:} Se modificó la regla \texttt{deconstruction\_expression} para permitir tanto declaraciones como asignaciones.

\subsection*{Atributo \texttt{AsyncMethodBuilder}}
\textbf{Problema:} No había una forma explícita de marcar métodos como constructores de métodos asíncronos. Sin embargo, desde C\# 10, se introdujo el atributo \texttt{AsyncMethodBuilder} para indicar métodos que construyen métodos asíncronos.

\textbf{Solución:} Se agregó soporte para el atributo \texttt{AsyncMethodBuilder}.

\subsection*{Atributo \texttt{CallerArgumentExpression}}
\textbf{Problema:} No había una forma explícita de obtener la expresión del argumento del llamador. Sin embargo, desde C\# 10, se introdujo el atributo \texttt{CallerArgumentExpression} para obtener esta información.

\textbf{Solución:} Se agregó soporte para el atributo \texttt{CallerArgumentExpression}.

\subsection*{Nuevo Formato para la Directiva \texttt{\#line}}
\textbf{Problema:} La directiva \texttt{\#line} no tenía un formato flexible. Sin embargo, desde C\# 10, se introdujo un nuevo formato para esta directiva que permite más flexibilidad.

\textbf{Solución:} Se modificó la regla \texttt{line\_directive} para soportar el nuevo formato.


\section*{C\# 11}

\subsection*{Literales de Cadena Raw}

\textbf{Problema:} Las cadenas literales no podían contener caracteres especiales sin escaparlos. Sin embargo, desde C\# 11, se introdujeron los literales de cadena raw, que permiten cadenas sin necesidad de escapar caracteres especiales.

\textbf{Solución:} Se agregó soporte para los literales de cadena raw en la regla \texttt{string\_literal}.

\subsection*{Soporte Genérico para Operaciones Matemáticas}

\textbf{Problema:} Las operaciones matemáticas no podían ser genéricas. Sin embargo, desde C\# 11, se introdujo el soporte genérico para operaciones matemáticas, permitiendo definir métodos que trabajen con tipos numéricos genéricos.

\textbf{Solución:} Se modificó la regla \texttt{type\_} para permitir tipos genéricos en operaciones matemáticas.

\subsection*{Atributos Genéricos}

\textbf{Problema:} Los atributos no podían ser genéricos. Sin embargo, desde C\# 11, se introdujo el soporte para atributos genéricos, permitiendo definir atributos que acepten tipos genéricos.

\textbf{Solución:} Se modificó la regla \texttt{attribute} para permitir atributos genéricos.

\subsection*{Literales de Cadena UTF-8}

\textbf{Problema:} Las cadenas literales no podían especificar codificación UTF-8 explícitamente. Sin embargo, desde C\# 11, se introdujo el soporte para literales de cadena UTF-8, permitiendo especificar la codificación explícitamente.

\textbf{Solución:} Se agregó soporte para literales de cadena UTF-8 en la regla \texttt{string\_literal}.

\subsection*{Saltos de Línea en Expresiones Interpoladas}

\textbf{Problema:} Las expresiones interpoladas no permitían saltos de línea dentro de la cadena. Sin embargo, desde C\# 11, se permite incluir saltos de línea en expresiones interpoladas.

\textbf{Solución:} Se modificó la regla \texttt{interpolated\_string\_expression} para permitir saltos de línea.

\subsection*{Patrones de Lista}

\textbf{Problema:} Los patrones no podían aplicarse a listas. Sin embargo, desde C\# 11, se introdujo el soporte para patrones de lista, permitiendo comparar listas con patrones específicos.

\textbf{Solución:} Se agregó soporte para patrones de lista en la regla \texttt{pattern}.

\subsection*{Tipos Locales a Archivo}

\textbf{Problema:} Los tipos definidos debían ser accesibles desde cualquier parte del archivo. Sin embargo, desde C\# 11, se introdujo el soporte para tipos locales a archivo, permitiendo definir tipos que solo sean accesibles dentro del mismo archivo.

\textbf{Solución:} Se modificó la regla \texttt{class\_definition} para permitir tipos locales a archivo.

\subsection*{Miembros Obligatorios}

\textbf{Problema:} No había una forma explícita de marcar miembros como obligatorios. Sin embargo, desde C\# 11, se introdujo el soporte para miembros obligatorios, permitiendo definir miembros que deben ser inicializados.

\textbf{Solución:} Se agregó soporte para miembros obligatorios en la regla \texttt{property\_declaration}.

\subsection*{Estructuras con Valores Predeterminados}

\textbf{Problema:} Las estructuras no tenían valores predeterminados. Sin embargo, desde C\# 11, se introdujo el soporte para estructuras con valores predeterminados, permitiendo que las estructuras tengan valores por defecto.

\textbf{Solución:} Se modificó la regla \texttt{struct\_declaration} para permitir valores predeterminados en estructuras.

\subsection*{Coincidencia de Patrones con \texttt{Span\textless char\textgreater}}

\textbf{Problema:} Los patrones no podían aplicarse a \texttt{Span\textless char\textgreater}. Sin embargo, desde C\# 11, se introdujo el soporte para coincidir patrones con \texttt{Span\textless char\textgreater}.

\textbf{Solución:} Se agregó soporte para patrones con \texttt{Span\textless char\textgreater} en la regla \texttt{pattern}.

\subsection*{Alcance Extendido de \texttt{nameof}}

\textbf{Problema:} El operador \texttt{nameof} solo podía aplicarse a nombres de tipos y miembros. Sin embargo, desde C\# 11, se extendió el alcance de \texttt{nameof} para permitir su uso en más contextos.

\textbf{Solución:} Se modificó la regla \texttt{expression} para permitir el uso extendido de \texttt{nameof}.

\subsection*{Tipo \texttt{IntPtr} Numérico}

\textbf{Problema:} El tipo \texttt{IntPtr} no era numérico. Sin embargo, desde C\# 11, se introdujo el soporte para que \texttt{IntPtr} sea un tipo numérico.

\textbf{Solución:} Se modificó la regla \texttt{type\_} para permitir que \texttt{IntPtr} sea un tipo numérico.

\subsection*{Campos \texttt{ref} y \texttt{scoped ref}}

\textbf{Problema:} Los campos no podían ser \texttt{ref}. Sin embargo, desde C\# 11, se introdujo el soporte para campos \texttt{ref} y \texttt{scoped ref}, permitiendo definir campos que se comporten como referencias.

\textbf{Solución:} Se agregó soporte para campos \texttt{ref} en la regla \texttt{field\_declaration}.

\subsection*{Mejora en la Conversión de Grupos de Métodos a Delegados}

\textbf{Problema:} La conversión de grupos de métodos a delegados no era tan flexible. Sin embargo, desde C\# 11, se mejoró esta conversión para permitir más flexibilidad en la asignación de métodos a delegados.

\textbf{Solución:} Se modificó la regla \texttt{expression} para mejorar la conversión de grupos de métodos a delegados.

\section*{C\# 12}

\subsection*{Constructores Primarios}

\textbf{Problema:} Los constructores deb\'ian definirse expl\'icitamente dentro de una clase o estructura. Desde C\# 12, se permite crear constructores primarios en cualquier tipo de clase o estructura.

\textbf{Soluci\'on:} Se modific\'o la regla \texttt{class\_definition} para permitir constructores primarios. Esto simplifica la definici\'on de clases y estructuras al combinar la declaraci\'on de par\'ametros con la inicializaci\'on de campos.

\subsection*{Expresiones de Colecci\'on}

\textbf{Problema:} Las colecciones se creaban utilizando el operador \texttt{new[]}. Desde C\# 12, se introdujo una nueva sintaxis para expresiones de colecci\'on, incluyendo el elemento de propagaci\'on (\texttt{..e}).

\textbf{Soluci\'on:} Se agreg\'o una nueva regla para soportar las expresiones de colecci\'on, definiendo la sintaxis de colecciones en l\'inea y el operador de propagaci\'on (\texttt{..e}).

\subsection*{Arrays en L\'inea}

\textbf{Problema:} Los arrays deb\'ian declararse con un tama\~no fijo utilizando llaves (\texttt{\{}). Desde C\# 12, se permite declarar arrays en l\'inea dentro de estructuras.

\textbf{Soluci\'on:} Se modific\'o la regla \texttt{type\_} para permitir arrays en l\'inea, simplificando la declaraci\'on de estructuras con arrays de tama\~no fijo.

\subsection*{Par\'ametros Opcionales en Expresiones Lambda}

\textbf{Problema:}Los par\'ametros en expresiones lambda no pod\'ian ser opcionales. Desde C\# 12, se permite el uso de par\'ametros opcionales en expresiones lambda.

\textbf{Soluci\'on:} Se actualiz\'o la regla \texttt{explicit\_anonymous\_function\_parameter} para soportar par\'ametros opcionales.

\subsection*{Par\'ametros \texttt{ref readonly}}

\textbf{Problema:} Los par\'ametros en expresiones lambda no pod\'ian ser \texttt{ref readonly}. Desde C\# 12, se permite el uso de par\'ametros \texttt{ref readonly} en expresiones lambda.

\textbf{Soluci\'on:} Se modific\'o la regla \texttt{explicit\_anonymous\_function\_parameter} para soportar par\'ametros \texttt{ref readonly}.

\subsection*{Alias Any Type}

\textbf{Problema:} El alias de tipo solo se pod\'ia utilizar con tipos nombrados. Desde C\# 12, se permite utilizar el alias de tipo con cualquier tipo.

\textbf{Soluci\'on:} Se actualiz\'o la regla \texttt{using\_directive} para permitir alias de cualquier tipo.

\subsection*{Atributo Experimental}

\textbf{Problema:} No hab\'ia una forma expl\'icita de marcar caracter\'isticas experimentales. Desde C\# 12, se introdujo el atributo \texttt{Experimental}.

\textbf{Soluci\'on:} Se agreg\'o soporte para el atributo \texttt{Experimental}.

